\section{Introduction}

Effectively parsing and understanding video content is a fundamental requirement for many downstream applications, such as video indexing, content retrieval, summarization, and automated editing. Within this pipeline, shot boundary detection (SBD) serves as a critical preliminary step. Its primary objective is to partition a video into temporally coherent shots by identifying the exact frames where camera views, scene contents, or editing states transition \cite{seridi2020survey,kar2024review}. While classical SBD methods rely on hand-crafted cues (e.g., color histograms and motion differences), these approaches have proven fragile under complex scenarios involving fast camera motion, flashes, and modern editing styles.

For SBD research, deep learning models have significantly advanced the state-of-the-art. Architectures like TransNet and TransNetV2 \cite{soucek2019transnet,soucek2020transnetv2} successfully model spatial-temporal representations directly from frame sequences. Concurrently, the rise of short-video platforms (e.g., TikTok, YouTube Shorts) has introduced a new, challenging operating regime \cite{vietnamnet2023shortform}. To address this, AutoShot \cite{zhu2023autoshot} targeted short-video SBD by leveraging neural architecture search (NAS) to find an optimal 3-D convolutional backbone, demonstrating strong performance on the newly introduced SHOT dataset. However, further improving these models encounters a severe bottleneck: the NAS pipeline requires immense computational overhead for supernet training and architecture evaluation. Furthermore, short-video SBD suffers from extreme class imbalance (boundary frames are exceptionally rare) and complex visual effects that act as hard negatives.

Following the promising strategy of separating representation learning from task-specific adaptation, we aim to study a more efficient scaling and refinement property for short-video SBD. Instead of re-running the expensive NAS process, we present a simple yet effective method to scale the performance of the existing architecture, termed \textbf{AutoShotV2}. Our core technical improvement is the introduction of a dual-branch classification head applied over the frozen 4864-dimensional features extracted by the AutoShot backbone. To handle the extreme class imbalance, we optimize this head using Focal Loss \cite{lin2017focal}, focusing the learning capacity on hard negatives and ambiguous transitions. Furthermore, we devise a boundary-aware many-hot supervision strategy in the auxiliary branch to provide denser temporal gradients without compromising the precise one-hot localization of the primary branch.

Finally, adapting deep models to threshold-based decision tasks often suffers from overconfidence and local noise. To mitigate this, we introduce a robust post-processing paradigm. At inference time, the primary boundary logits are calibrated using temperature scaling \cite{guo2017calibration} to adjust the numerical scale of the predictions without altering their ranking. Subsequently, a one-dimensional Gaussian temporal filter is applied to smooth out isolated, false-positive peaks caused by sudden flashes or camera shakes.

Based on these core designs---frozen robust representations, dual-branch focal supervision, and calibrated temporal smoothing---we develop AutoShotV2. On the SHOT test set, AutoShotV2 achieves an F1-score of \PaperASVTwoShotFOne, the best result among the evaluated methods, improving over the reproduced AutoShot baseline and the reported TransNetV2 baseline by \PaperASVTwoVsAutoShotPP{} and \PaperASVTwoVsTransNetPP{} F1 points, respectively.

In summary, our main contributions are:
\begin{itemize}
\item We propose AutoShotV2, an efficient and highly accurate SBD pipeline that significantly improves short-video boundary detection without the computational burden of retraining the NAS-searched backbone.
\item We design a novel two-branch classification head trained with Focal Loss and boundary-aware many-hot supervision to address the extreme class imbalance inherent in short videos.
\item We introduce a calibrated inference pipeline leveraging temperature scaling and Gaussian smoothing to ensure stable, threshold-based boundary decisions.
\item We provide extensive experiments on SHOT, BBC, and ClipShots, demonstrating strong performance gains and identifying important cross-domain limitations for future research.
\end{itemize}

